\documentclass[aspectratio=169,11pt]{beamer}
\usepackage[utf8]{inputenc}
\usepackage[spanish]{babel}
\usepackage{amsmath}
\usepackage{graphicx}
\usepackage{xcolor}
\usepackage{tikz}

\definecolor{ejercicio}{RGB}{59, 88, 102}
\definecolor{acento}{RGB}{200, 80, 50}

\usetheme{default}
\usecolortheme{default}
\setbeamercolor{title}{fg=ejercicio}
\setbeamercolor{frametitle}{fg=ejercicio}
\setbeamercolor{structure}{fg=ejercicio}
\setbeamertemplate{navigation symbols}{}
\setbeamertemplate{footline}[frame number]

\newcommand{\foto}[2]{%
	\IfFileExists{#1}{\includegraphics[width=#2]{#1}}{%
		\fbox{\parbox[c][3.2cm][c]{#2}{\centering\small\color{gray}[ foto: \detokenize{#1} ]}}}}

\newcommand{\dato}[2]{%
	\begin{minipage}{0.3\textwidth}\centering
		{\Huge\bfseries\color{acento} #1}\\[4pt]
		{\small #2}
\end{minipage}}

\title{Design Review}
\subtitle{Proyectos III. Entrega 1}
\author{César Vidal Fernández}
\date{Septiembre de 2026}

\begin{document}
	
	\begin{frame}
		\titlepage
	\end{frame}
	
	% ---------------------------------------------------------------
	\begin{frame}{El punto de partida: el Bernabéu}
		\vspace{0.4cm}
		\begin{center}
			\dato{225 M\texteuro}{de inversión en el césped retráctil}
			\hfill
			\dato{30 m}{de profundidad del foso (\textit{hipogea})}
			\hfill
			\dato{15 min}{en desplegar la cubierta}
		\end{center}
		\vspace{0.6cm}
		\begin{center}
			\foto{imagen/bernabeuViejo.jpg}{0.4\textwidth}
			\hspace{0.5cm}
			\foto{imagen/bernabeuNuevo.png}{0.4\textwidth}
		\end{center}
	\end{frame}
	
	% ---------------------------------------------------------------
	\begin{frame}{Por qué me centro en otra cosa}
		\vspace{0.5cm}
		\begin{center}
			{\Large Motor de físicas adaptable\\}
		\end{center}
		\vspace{0.6cm}
		\begin{center}
			\begin{tabular}{ccc}
				\textbf{TFG} & $\longrightarrow$ & \textbf{Proyectos III} \\[4pt]
				Audio y estadio & & Avión rompiendo\\ & & la barrera del sonido
			\end{tabular}
		\end{center}
	\end{frame}
	
	% ---------------------------------------------------------------
	\begin{frame}{Objetivo}
		\vspace{0.4cm}
		\begin{center}
			{\Large Simulación \textbf{en tiempo real}\\ del instante en que se supera Mach 1.}
		\end{center}
		\vspace{0.6cm}
		\begin{center}
			\foto{imagen/concorde.jpg}{0.5\textwidth}
		\end{center}
	\end{frame}
	
	% ---------------------------------------------------------------
	\begin{frame}{Algunos datos}
		\vspace{0.4cm}
		\begin{center}
			\dato{343 m/s}{velocidad del sonido a 20\,$^\circ$C}
			\hfill
			\dato{1\,235 km/h}{lo mismo en km/h}
			\hfill
			\dato{1947}{primer vuelo supersónico (Bell X-1)}
		\end{center}
		\vspace{0.8cm}
		\begin{center}
			{\large A mayor altura hace más frío, y el sonido va más lento:\\
				Mach 1 son unos \textbf{295 m/s} a 11\,000 m.}
		\end{center}
	\end{frame}
	
	% ---------------------------------------------------------------
	\begin{frame}{El cono de Mach}
		\begin{columns}[c]
			\column{0.55\textwidth}
			\begin{center}
				\begin{tikzpicture}[scale=0.85]
					% frentes de onda emitidos (Mach 2)
					\foreach \d in {0.8,1.6,2.4,3.2}{
						\draw[ejercicio!60] (-\d,0) circle (\d/2);
					}
					% cono
					\draw[acento, thick] (0,0) -- ({-4.6*cos(30)},{4.6*sin(30)});
					\draw[acento, thick] (0,0) -- ({-4.6*cos(30)},{-4.6*sin(30)});
					% avión
					\fill[black] (0,0) circle (2.5pt);
					\draw[->, thick] (0.2,0) -- (1.2,0);
					\node[right] at (1.2,0) {\small avión};
					\node[acento] at (-1.2,0.55) {\small $\mu$};
				\end{tikzpicture}
			\end{center}
			\column{0.4\textwidth}
			\begin{center}
				{\Large $\displaystyle \mu = \arcsin\left(\frac{1}{M}\right)$}
			\end{center}
		\end{columns}
	\end{frame}
	
	% ---------------------------------------------------------------
	\begin{frame}{Cuanto más rápido, más cerrado}
		\vspace{0.5cm}
		\begin{center}
			\renewcommand{\arraystretch}{1.6}
			\begin{tabular}{c|cccc}
				\textbf{Mach} & 1,2 & 1,5 & 2 & 3 \\ \hline
				\textbf{Ángulo del cono} & 56,4$^\circ$ & 41,8$^\circ$ & 30$^\circ$ & 19,5$^\circ$
			\end{tabular}
		\end{center}
	\end{frame}
	
	% ---------------------------------------------------------------
	\begin{frame}{Qué se podrá tocar}
		\begin{columns}[t]
			\column{0.3\textwidth}
			\begin{center}
				{\large\bfseries\color{ejercicio} Entrada}\\[6pt]
				Temperatura\\
				Altitud\\
				Velocidad inicial\\
				Forma del avión\\
				Observador
			\end{center}
			\column{0.3\textwidth}
			\begin{center}
				{\large\bfseries\color{ejercicio} Variables en ejecución}\\[6pt]
				Aceleración\\
				Velocidad objetivo
			\end{center}
			\column{0.3\textwidth}
			\begin{center}
				{\large\bfseries\color{ejercicio} Salida}\\[6pt]
				Número de Mach\\
				Cono de Mach\\
				Campo de presión\\
				Estampido\\
				FPS
			\end{center}
		\end{columns}
	\end{frame}
	
	% ---------------------------------------------------------------
	\begin{frame}{Simplificaciones}
		\vspace{0.4cm}
		\begin{center}
			{\Large
				Vuelo recto y a altura constante\\[6pt]
				Aire ideal, con viento y turbulencias pero sin viscosidad\\[6pt]
				Avión como objeto simple (geometría reducida)\\[6pt]
				Claridad antes que exactitud}
		\end{center}
	\end{frame}
	
	% ---------------------------------------------------------------
	\begin{frame}
		\vspace{1.2cm}
		\begin{center}
			{\Huge\color{ejercicio} Gracias.}
		\end{center}
	\end{frame}
	
\end{document}